\SourceLeaf{60}
$(1+n\cos\varphi)^{\nu-1}$ coefficientes derivari queant. Sic enim ab ipsa $(1+n\cos\varphi)^{-1}$ transire licebit ad potestates negativas altiores, vitatis seriebus pro singulo coefficiente infinitis.

$\alpha$) Ad scrutandam, cui subiecta est, posito $\nu=-1$, coefficiens initialis $A$, legem, in expressione indefinita termini eiusdem $h^{\text{ti}}$ ponamus $\nu=-1$, provenietque, ob
\[
{}^{-1}\mathfrak{B}^{(2h)} = 1, \qquad \text{ille} \qquad = \frac{1}{2^{2h}}\,{}^{2h}\mathfrak{B}^{(h)}\,n^{2h}.
\]
Est autem
\[
\frac{1}{2^{2h}}\,{}^{2h}\mathfrak{B}^{(h)} = \frac{1}{2^{2h}}\,\frac{2h.(2h-1).\,.\,.\,.(h+1)}{1.2.\,.\,.\,. h}
\]
\[
= \frac{1}{2^{2h}} . \frac{2h.(2h-1).\,.\,.\,. 1}{1.2.\,.\,.\,.\,. h} . \frac{1}{h.(h-1).\,.\,. 1}
\]
\[
= \frac{1}{2^{2h}}\,\frac{(2h-1)(2h-3).\,.\,.\,. 1}{1.2.\,.\,. h} . \frac{2h.(2h-2).\,.\,. 2}{h.(h-1).\,.\,.\,.\,. 1}
\]
\[
= \frac{1}{2^{h}} . \frac{(2h-1)(2h-3).\,.\,. 1}{1.2.\,.\,.\,. h},
\]
fitque his transformationibus terminus ille $h^{\text{tus}}$
\[
= \frac{1}{2^{h}}\,\frac{(2h-1).(2h-3).\,.\,. 1}{1.2.\,.\,.\,. h} .  n^{2h}.
\]
Quare, cum haec expressio sit manifesto quantitatis $(1-nn)^{-\frac{1}{2}}$, in seriem evolutae, terminus $h^{\text{tus}}$, sequitur, fore ipsam $A$, casu quo $\nu=-1$,
\[
= (1-nn)^{-\frac{1}{2}}
\]
\par\noindent\hspace*{\parindent}$\beta$) Coefficiens post initialem primus $A^{(1)}$ definitur relatione ante demonstrata (7, $\delta$)
\[
A^{(1)} = \frac{2.(\nu+2)}{n}\int An\,\mathrm{d}n - 2n.A,
\]
quae hic, ob $\nu=-1$, mutatur in
\[
A^{(1)} = \frac{2}{n}\int An\,\mathrm{d}n - 2nA = \frac{2}{n}\int(1-nn)^{-\frac{1}{2}}n\,\mathrm{d}n - 2n(1-nn)^{-\frac{1}{2}}.
\]
Peracta integratione (sic instituenda, ut cum $n$ et integrale evanescat) habebitur
\[
A^{(1)} = -\frac{2}{n}(1-nn)^{\frac{1}{2}} + \frac{2}{n} - 2n(1-nn)^{-\frac{1}{2}},
\]
sive brevius
\[
A^{(1)} = \frac{2}{n}\bigl(1-(1-nn)^{-\frac{1}{2}}\bigr).
\]

